% An invented source paper for the Sextant demo. Nothing in it is real.
\documentclass[10pt]{article}
\usepackage[T1]{fontenc}
\usepackage[margin=2.4cm]{geometry}
\usepackage{lmodern}
\pagestyle{empty}
\begin{document}
\begin{center}
{\small\itshape Proceedings of the Fictional Society for Urban Ecology, pp.~57--64, 2022\par}\vspace{10pt}
{\LARGE\bfseries Compost Amendments and Water Infiltration in Compacted Urban Soils\par}\vspace{8pt}
{\large Riley Testwell and Casey Mockford\par}\vspace{4pt}
{\small Department of Imaginary Studies, Nowhere University\par}
\end{center}
\vspace{6pt}
\noindent\textbf{Abstract.} Compacted urban soils absorb little of the rain that falls on them. We tested whether a single compost amendment changes how much simulated rainfall a plot takes in.

\section*{Methods}
Eighteen plots on a former car park were loosened by hand; half received a 5 cm layer of compost worked into the top layer. A rainfall simulator delivered a steady shower to each plot, and run-off was collected and measured.

\section*{Results}
Averaged across the nine treated plots and measured from the moment the shower began, compacted plots with compost amendments absorbed 31 mm of simulated rainfall in the first hour, against 9 mm on untreated plots. The difference persisted when the test was repeated the following spring.

\section*{Discussion}
Compost improves the structure of the soil rather than its surface, so its effect lasts beyond the season it is applied. It pairs well with a mulch, which keeps the absorbed water from evaporating.
\end{document}
